% Chapter III, Part 2: Sections 29-35 + Bibliographical Notes
% The Thermal Instability of a Layer of Fluid Heated from Below
% 2. THE EFFECT OF ROTATION (continued)
% Equations continue from eq:3-210 in chapter_3_part1.tex

%% --- Section 29 ---
\section{On the onset of convection as overstability. The solution for the case of two free boundaries}\label{sec:3-29}

We now take up the question, set aside in \S~26, whether or not instability can arise as oscillations of increasing amplitude, i.e., as overstability. This requires us to return to the general equations~\eqref{eq:3-93}--\eqref{eq:3-95} which include the time constant~$\sigma$. In \S~30 we shall describe how best one may treat these equations to decide under what conditions stationary or overstable convection arise. In this section we shall restrict ourselves to the case of two free boundaries; for in this case the problem can be solved by elementary methods and it suggests how one may treat the general case.

By applying the operator $(D^2 - a^2 - \sigma)(D^2 - a^2 - p\sigma)$ to equation~\eqref{eq:3-95}, we can eliminate $Z$ and $\Theta$ and obtain
\begin{equation}\label{eq:3-211}
(D^2 - a^2 - p\sigma)(D^2 - a^2 - \sigma)^2(D^2 - a^2) + \mathscr{T} D^2] W = -Ra^2(D^2 - a^2 - \sigma)W.
\end{equation}

As in \S~27\,$(a)$, we can show that in this case also the proper solution for $W$ belonging to the lowest mode is
\begin{equation}\label{eq:3-212}
W = W_0 \sin \pi z.
\end{equation}

Substituting this solution for $W$ in equation~\eqref{eq:3-211}, we obtain the characteristic equation
\begin{equation}\label{eq:3-213}
(\pi^2 + a^2 + p\sigma)[(\pi^2 + a^2 + \sigma)^2(\pi^2 + a^2) + \pi^2\mathscr{T}] = Ra^2(\pi^2 + a^2 + \sigma),
\end{equation}
where it must be remembered that $\sigma$ can be complex. Letting
\begin{equation}\label{eq:3-214}
x = \frac{a^2}{\pi^2}, \quad i\sigma_1 = \frac{\sigma}{\pi^2}, \quad R_1 = \frac{R}{\pi^4}, \quad \text{and} \quad T_1 = \frac{\mathscr{T}}{\pi^4},
\end{equation}
we can rewrite equation~\eqref{eq:3-213} in the form
\begin{equation}\label{eq:3-215}
R_1 = \frac{1}{x}(1 + x + ip\sigma_1)(1 + x)(1 + x + i\sigma_1) + \frac{T_1}{1 + x + i\sigma_1}\bigg\}.
\end{equation}

It is apparent from equation~\eqref{eq:3-215} that for an arbitrarily assigned $\sigma_1$, $R_1$ will be complex. But the physical meaning of $R_1$ requires it to be real. Consequently, the condition that $R_1$ be real implies a relation between the real and imaginary parts of $\sigma_1$. Since our principal interest is to specify the critical Rayleigh number for the onset of instability via a state of purely oscillatory motions, we shall in the first instance suppose that $\sigma_1$ in equation~\eqref{eq:3-215} is real and seek the conditions for such solutions to exist. This will suffice to answer the principal question as to when instability will set in as stationary convection and when as overstable oscillations. To establish that the criteria so obtained are both necessary and sufficient, we must consider the characteristic equation~\eqref{eq:3-215} more generally, and this we will do in \S~$(a)$ below.

We shall assume, then, that in equation~\eqref{eq:3-215} $\sigma_1$ is real. On collecting the real and imaginary parts of the expression on the right-hand side of equation~\eqref{eq:3-215}, we have
\begin{equation}\label{eq:3-216}
R_1 = \frac{1+x}{x}\bigg\{(1+x)^2 - p\sigma_1^2 + \frac{T_1}{1+x}\,\frac{(1+x)^2 + p\sigma_1^2}{(1+x)^2 + \sigma_1^2} + i\sigma_1\bigg[(1+x)(1+p) - T_1\,\frac{(1-p)}{(1+x)^2 + \sigma_1^2}\bigg]\bigg\}.
\end{equation}

The real and the imaginary parts of this equation must vanish separately. We thus obtain the pair of equations:
\begin{equation}\label{eq:3-217}
R_1 = \frac{1+x}{x}\bigg\{(1+x)^2 - p\sigma_1^2 + \frac{T_1}{1+x}\,\frac{(1+x)^2 + p\sigma_1^2}{(1+x)^2 + \sigma_1^2}\bigg\}
\end{equation}
and
\begin{equation}\label{eq:3-218}
(1+x)(1+p) = T_1\,\frac{(1-p)}{(1+x)^2 + \sigma_1^2}.
\end{equation}

A relation which follows from equation~\eqref{eq:3-218} is
\begin{equation}\label{eq:3-219}
\frac{T_1}{1+x}\,\frac{(1+x)^2 + p\sigma_1^2}{(1+x)^2 + \sigma_1^2} = \frac{T_1}{1+x}\,\frac{T_1(1-p)\sigma_1^2}{(1+x)^2 + \sigma_1^2} = \frac{T_1}{1+x} - (1+p)\sigma_1^2.
\end{equation}

Using this relation in~\eqref{eq:3-217}, we have
\begin{equation}\label{eq:3-220}
R_1 = \frac{1}{x}\big[(1+x)^3 + T_1 - (1+x)(1+2p)\sigma_1^2\big].
\end{equation}

Next, solving equation~\eqref{eq:3-218} for $\sigma_1^2$, we find
\begin{equation}\label{eq:3-221}
\sigma_1^2 = \frac{T_1}{1+x}\,\frac{1-p}{1+p} - (1+x)^2.
\end{equation}

Using this expression for $\sigma_1^2$ in~\eqref{eq:3-220}, we obtain on further simplification the result
\begin{equation}\label{eq:3-222}
R_1 = 2(1+p)\frac{1}{x}\bigg[(1+x)^3 + \frac{p^2}{(1+p)^2}\,T_1\bigg].
\end{equation}

Equations~\eqref{eq:3-221} and~\eqref{eq:3-222} are the equations which must be satisfied if overstability is to occur for a wave number corresponding to $x$ and a Taylor number corresponding to $T_1$.

One conclusion which equation~\eqref{eq:3-221} enables us to draw, at once, is that solutions describing overstability cannot occur if
\begin{equation}\label{eq:3-223}
\frac{T_1}{1+x}\,\frac{1-p}{1+p} < 1;
\end{equation}
for $\sigma_1^2$ would then be negative, contrary to hypothesis. Clearly, \eqref{eq:3-223} will, \emph{a fortiori}, be the case if $p > 1$. Accordingly, \emph{for $p > 1$, overstability cannot occur and the principle of the exchange of stabilities is valid}.

For overstability to be at all possible, $p$ must be less than one; and even when this is the case, we shall obtain real frequencies, $\sigma_1$, only if
\begin{equation}\label{eq:3-224}
T_1 > \frac{1+p}{1-p}(1+x)^3.
\end{equation}

For a given $T_1$, overstable solutions are therefore possible only for $x < x_a$ where $x_a$ is such that
\begin{equation}\label{eq:3-225}
(1+x_a)^3 = T_1\,\frac{1-p}{1+p}.
\end{equation}

When $x = x_a$ and $\sigma_1^2 = 0$ and (cf.\ equation~\eqref{eq:3-220})
\begin{equation}\label{eq:3-226}
R_1 = \frac{1}{x_a}\big[(1+x_a)^3 + T_1\big];
\end{equation}
and this is, as one should expect, the value of $R_1$ at which stationary convection will occur for a wave number corresponding to $x_a$ (see equation~\eqref{eq:3-130}). For $x > x_a$, overstability is not possible for the given $p$ and $T_1$; and the onset of instability as stationary convection remains the only possibility. For $x < x_a$, overstability is possible; and the question of discriminating between the two manners of instability arises. In making the discrimination, we shall suppose, other things being equal, that manner of instability will occur that allows a solution for a lower Rayleigh number. While this is intuitively obvious, it requires proof; this is furnished in \S~$(a)$ below.

Consider the $(x, R_1)$-plane (see Fig.~27). In this plane, we have first the curve,
\begin{equation}\label{eq:3-227}
R_1^{(c)} = \frac{1}{x}\big[(1+x)^3 + T_1\big],
\end{equation}
which defines the locus of states which are marginal with respect to stationary convection. The overstable solutions branch off from this

\figurePlaceholder{27}{The disposition of the curves for marginal stability in the $(x, R_1)$-plane for a Taylor number, $T = 10^6$, and for various values of the Prandtl number~$p$. The curve labelled `convection' applies for all values of~$p$ and defines the locus $R_1^{(c)}$ (equation~227). The remaining curves are the overstable loci $R_1^{(o)}$ for values of~$p$ by which they are labelled. The minimum for the curve $C$ ($p = 0{\cdot}6120$) and for the convection curve occurs for the same value of~$R_1$.}

\noindent locus at the point $x_a$ (see equation~\eqref{eq:3-226}); and for $x < x_a$ they are described by (cf.\ equations~\eqref{eq:3-220} and~\eqref{eq:3-222})
\begin{equation}\label{eq:3-228}
R_1^{(o)} = R_1^{(c)} - (1+2p)\sigma_1^2\,\frac{1+x}{x} = 2(1+p)\frac{1}{x}\bigg[(1+x)^3 + \frac{p^2}{(1+p)^2}\,T_1\bigg],
\end{equation}
where the first form shows that $R_1^{(o)}$ (when this branch of the solution exists) is \emph{always} less than $R_1^{(c)}$. Depending on $p$ and $T_1$, the branch point $x_a$ can occur either before, or after, the point, $x_{\min}^{(c)}$, at which $R_1^{(c)}$ attains its minimum. If $x_a > x_{\min}^{(c)}$, then it is clear that for all $x < x_a$, overstability is the preferred manner of instability: these are shown by the different curves labelled by $B$, $C$, $D$, and $E$ in Fig.~27. It is apparent from this figure that the case which will distinguish whether, with increasing $R_1$ (for a given $T_1$), overstability or stationary convection will first manifest itself, is the one for which the minima of the two curves $R_1^{(c)}$ and $R_1^{(o)}$ are equal. Thus, if the branch point $x_a$ occurs for a value of $x$ less than what corresponds to $C$ in Fig.~27, we can exclude overstability; on the other hand, if the branch point occurs for a value of $x$ greater than what corresponds to $C$, we can exclude stationary convection.

We shall now show that there exists a value of $p$ such that $R_{1,\min}^{(c)}$ approaches $R_{1,\min}^{(o)}$ from above as $T_1 \to \infty$. For according to equations~\eqref{eq:3-227} and~\eqref{eq:3-228}, the asymptotic behaviours of $R_{1,\min}^{(c)}$ and $R_{1,\min}^{(o)}$ are given by (cf.\ equation~\eqref{eq:3-139})
\begin{equation}\label{eq:3-229}
R_{1,\min}^{(c)} \to 3(\tfrac{4}{3}T_1)^{\frac{2}{3}},
\end{equation}
\begin{equation}\label{eq:3-230}
R_{1,\min}^{(o)} \to 2(1+p)\bigg\{\tfrac{3}{2}\bigg[\frac{p^2}{2\,(1+p)^2}\,T_1\bigg]^{\frac{2}{3}}\bigg\}.
\end{equation}

The condition that $R_{1,\min}^{(c)} \to R_{1,\min}^{(o)}$ as $T_1 \to \infty$, clearly, requires that
\begin{equation}\label{eq:3-231}
2(1+p)\bigg[\frac{p^2}{(1+p)^2}\bigg]^{\frac{2}{3}} = 1,
\end{equation}
or
\begin{equation}\label{eq:3-232}
\frac{p^{\frac{4}{3}}}{(1+p)^{\frac{1}{3}}} = 1.
\end{equation}

It is found that the required root of this equation is
\begin{equation}\label{eq:3-233}
p = 0{\cdot}67659 = p^* \;\text{(say)}.
\end{equation}

From the monotone dependence of the minimum of the curve,
\begin{equation}\label{eq:3-234a}
y = \frac{1}{x}\big[(1+x)^3 + X\big],
\end{equation}
on $X$, we conclude that for $p > p^*$, $R_{1,\min}^{(c)} < R_{1,\min}^{(o)}$ for all $T_1$. Hence, for $1 > p > p^*$, the situations indicated by $D$ or $E$ will prevail for all $T_1$. Therefore, \emph{for $p > p^*$, instability will always manifest itself, first, as stationary convection}.

For $p < p^*$, the situation indicated by $C$ occurs for a finite $T_1$, say $T_1^{(p)}$. For $T_1 < T_1^{(p)}$, the disposition of the curves $R_1^{(o)}$, with respect to $R_1^{(c)}$, is as shown by the curves $D$ and $E$ in Fig.~27; when $T_1 = T_1^{(p)}$, the situation is as shown by the curve $C$; and for $T_1 > T_1^{(p)}$ the disposition of the curves $R_1^{(o)}$, with respect to $R_1^{(c)}$, is as shown by the curves $B$ and $A$. We conclude, then, that \emph{for $p < p^*$ there exists a $T_1^{(p)}$ such that for $T_1 < T_1^{(p)}$, the onset of instability will be as stationary convection, while for $T_1 > T_1^{(p)}$ it will be as overstability}.

There is no simple formula which gives $T_1^{(p)}$ as a function of $p$: it is simply determined by the condition that $R_{1,\min}^{(c)}$ and $R_{1,\min}^{(o)}$ are equal for

\begin{table}[ht]
\centering
\caption{Taylor numbers above which instability occurs as overstability}
\label{tab:X}
\begin{tabular}{cccccc}
\hline
$p$ & $T^{(p)}$ & $R$ & $p$ & $T^{(p)}$ & $R$ \\
\hline
$0$ & & & $0{\cdot}55$ & $15{,}370$ & 7748 \\
$0{\cdot}1$ & 548 & 1315 & $0{\cdot}60$ & $68{,}150$ & $16{,}290$ \\
$0{\cdot}2$ & 728 & 1431 & $0{\cdot}63$ & $2{\cdot}588\times 10^5$ & $3{\cdot}870\times 10^4$ \\
$0{\cdot}3$ & 990 & 1669 & $0{\cdot}65$ & $1{\cdot}233\times 10^6$ & $1{\cdot}059\times 10^5$ \\
$0{\cdot}4$ & 3103 & 2890 & $0{\cdot}6766$ & $\infty$ & $\infty$ \\
$0{\cdot}5$ & 8595 & 4970 & & & \\
\hline
\end{tabular}
\end{table}

\figurePlaceholder{28}{The variation of $T^{(p)}$ with the Prandtl number~$p$.}

\noindent $T_1 = T_1^{(p)}$. Values of $T^{(p)}$ determined by this condition are given in Table~X for a few values of $p$; and Fig.~28 shows the division of the $(p, T)$-plane into regions where the manner of instability at onset is one or the other; the part of this plane where overstability can occur for Rayleigh numbers higher than are required for the onset of instability is bounded by $p = 1$ and the locus~$T^{(p)}$.

For $p < p^*$, the complete $(R_c, T)$-relation can be deduced from the results given in Table~VII by a simple transformation of scales: for, according to equation~\eqref{eq:3-222}, we have only to interpret $T$ in Table~VII

\begin{table}[ht]
\centering
\caption{Critical Rayleigh numbers and related constants for the onset of overstability in case both bounding surfaces are free and $p = 0{\cdot}025$}
\label{tab:XI}
\small
\begin{tabular}{ccccccc}
\hline
$T$ & $a_c$ & $R_c$ & $p|\sigma|$ & & & \\
\hline
0 & $2{\cdot}33$ & & $1{\cdot}348\times 10^4$ & & \text{imaginary} & \\
$1{\cdot}681\times 10^3$ & $2{\cdot}70$ & $1{\cdot}014\times 10^4$ & $1{\cdot}388\times 10^5$ & & $1{\cdot}504$ & \\
$1{\cdot}681\times 10^4$ & $2{\cdot}944$ & $3{\cdot}079\times 10^4$ & $1{\cdot}064\times 10^6$ & & $1{\cdot}262$ & \\
$8{\cdot}493\times 10^4$ & $3{\cdot}278$ & $6{\cdot}184\times 10^4$ & $2{\cdot}615\times 10^6$ & & $1{\cdot}349$ & \\
$1{\cdot}875\times 10^5$ & $3{\cdot}710$ & $8{\cdot}055\times 10^4$ & $3{\cdot}435\times 10^6$ & & $1{\cdot}266$ & \\
$3{\cdot}364\times 10^5$ & $4{\cdot}220$ & $1{\cdot}007\times 10^5$ & $4{\cdot}713\times 10^6$ & & $1{\cdot}164$ & \\
$8{\cdot}495\times 10^5$ & $5{\cdot}011$ & $1{\cdot}303\times 10^5$ & $7{\cdot}513\times 10^6$ & & $1{\cdot}037$ & \\
$1{\cdot}681\times 10^6$ & $5{\cdot}968$ & $1{\cdot}937\times 10^5$ & $1{\cdot}183\times 10^7$ & & $0{\cdot}944$ & \\
$3{\cdot}943\times 10^6$ & $6{\cdot}961$ & $2{\cdot}851\times 10^5$ & $2{\cdot}092\times 10^7$ & & $0{\cdot}5029$ & \\
$1{\cdot}681\times 10^7$ & $8{\cdot}606$ & $4{\cdot}330\times 10^5$ & $4{\cdot}858\times 10^7$ & & $0{\cdot}5608$ & \\
$3{\cdot}943\times 10^7$ & $10{\cdot}43$ & $6{\cdot}360\times 10^5$ & $7{\cdot}728\times 10^7$ & & $0{\cdot}5618$ & \\
$1{\cdot}681\times 10^8$ & $12{\cdot}86$ & $9{\cdot}497\times 10^5$ & $1{\cdot}501\times 10^8$ & & $0{\cdot}4931$ & \\
$3{\cdot}943\times 10^8$ & $13{\cdot}06$ & $1{\cdot}370\times 10^6$ & $1{\cdot}784\times 10^8$ & & $0{\cdot}4790$ & \\
$1{\cdot}681\times 10^9$ & $17{\cdot}05$ & $2{\cdot}067\times 10^6$ & $3{\cdot}447\times 10^8$ & & $0{\cdot}3896$ & \\
$1{\cdot}681\times 10^{10}$ & $28{\cdot}02$ & $4{\cdot}470\times 10^6$ & $8{\cdot}958\times 10^8$ & & $0{\cdot}2175$ & \\
$1{\cdot}681\times 10^{11}$ & $38{\cdot}87$ & $8{\cdot}470\times 10^6$ & $8{\cdot}476\times 10^9$ & & $0{\cdot}0320$ & \\
$1{\cdot}681\times 10^{12}$ & $191{\cdot}51$ & $2{\cdot}075\times 10^7$ & & & & \\
\hline
\end{tabular}
\end{table}

\noindent as now signifying $p^4T/(1+p)^2$ and multiply the values under $R_c$ by $2(1+p)$. The $(R_c, T)$-relations derived in this manner are illustrated in Fig.~29. For later comparisons with results derived for other boundary conditions, the numerical form of the relation for $p = 0{\cdot}025$ is given in Table~XI.

It will be observed that in accordance with the result given in~\eqref{eq:3-229}, along the various overstable solutions $R \to \text{constant}\cdot T^{\frac{2}{3}}$ as $T \to \infty$. The explicit forms of the asymptotic behaviours of the Rayleigh number and the wave number of the associated disturbance for the onset of overstability may be noted; they are:
\begin{equation}\label{eq:3-234}
R_c^{(o)} \to \frac{6p^{\frac{4}{3}}}{(1+p)^{\frac{1}{3}}}\,(\tfrac{1}{2}\pi^2 T)^{\frac{2}{3}} \quad (p^2 T \to \infty).
\end{equation}

\begin{equation}\label{eq:3-235}
a_c^{(o)} \to \bigg(\frac{p}{1+p}\bigg)^{\frac{1}{3}} (\tfrac{1}{2}\pi^2 T)^{\frac{1}{6}} \quad (p^2 T \to \infty).
\end{equation}

The corresponding behaviour of $|\sigma| = \pi^2\sigma_1$ can be deduced from equation~\eqref{eq:3-221}. We find
\begin{equation}\label{eq:3-236}
|\sigma| \to \frac{(2-3p)^{\frac{1}{2}}}{|p(1+p)|^{\frac{1}{2}}}\,(\tfrac{1}{2}\pi^2 T)^{\frac{1}{2}} \quad (p^2 T \to \infty),
\end{equation}
where it may be noted that according to the definition of $\sigma$ (see equation~\eqref{eq:3-92}),
\begin{equation}\label{eq:3-237}
|p| = \frac{2|\sigma|}{\Omega} = \frac{2|\sigma_1|}{\sqrt{T_1}}.
\end{equation}

A further relation which follows from the foregoing is
\begin{equation}\label{eq:3-238}
|p|a_c^{(o)} \to 2\pi\Omega^{\frac{1}{2}}\,\frac{(1-1{\cdot}5p^{\frac{1}{2}})}{1+p} \quad (p^2 T \to \infty).
\end{equation}

It should be noted that the asymptotic relations~\eqref{eq:3-234}, \eqref{eq:3-236}, and~\eqref{eq:3-238} are valid only for $p^2 T \to \infty$, \emph{not} for $T \to \infty$, a distinction which is important when $p \to 0$ (see \S~32).

\figurePlaceholder{29}{The $(R_c, T)$-relations for a rotating horizontal layer of fluid heated below. The curves have been derived for the case when both bounding surfaces are free. The curve labelled `convection curve' is the $(R_c, T)$-relation for the onset of ordinary cellular convection. The remaining curves are the corresponding relations for the onset of overstability. The values of~$p$ to which the various curves refer are shown at the top of each curve. It will be seen that for each value of $p < 0{\cdot}677$, the instability sets in as ordinary cellular convection for $T$ less than a certain $T^{(p)}$ while it sets in as overstability for $T > T^{(p)}$.}

\subsection*{(a) The nature of the roots of the characteristic equation~\eqref{eq:3-215}}

The preceding discussion of the dependence of the manner of instability on $p$ and $T$ was carried out in terms of the explicit solutions for two special cases of equation~\eqref{eq:3-215}: the case when $\sigma_1 = 0$ and the marginal state is a stationary one; and the case when $\sigma_1$ is real and the marginal state is an oscillatory one. It was further assumed that the manner of instability which will occur at onset is the one which allows a solution for a lower Rayleigh number. It was pointed out that this assumption requires justification; and this we shall now provide.

In considering the roots of equation~\eqref{eq:3-215}, generally, we shall find it convenient to replace $i\sigma_1$ by $\sigma_1$ so that overstability now corresponds

\noindent to $\sigma_1$ being a pure imaginary number. With this replacement, the equation we have to consider may be written as
\begin{equation}\label{eq:3-239}
(1+x+p\sigma_1)\big[(1+x+\sigma_1)^2(1+x) + \frac{T_1}{1+x}\,d^2\big] = R_1(1+x+\sigma_1)\,\frac{x}{1+x},
\end{equation}
where
\begin{equation}\label{eq:3-240}
\begin{aligned}
B &= \frac{1}{p}(1+x)(1+2p), \\[4pt]
C &= \frac{1}{p}\big[(2+p)(1+x)^2 + p\,\frac{T_1}{1+x} - R_1\,\frac{x}{1+x}\big], \\[4pt]
D &= \frac{1}{p}\big[(1+x)^3 + T_1 - R_1\,x\big].
\end{aligned}
\end{equation}

This is a cubic equation for $\sigma_1$; its explicit form is
\begin{equation}\label{eq:3-241}
\sigma_1^3 + B\sigma_1^2 + C\sigma_1 + D = 0,
\end{equation}

A combination of the coefficients of the cubic~\eqref{eq:3-241}), which is important for the characterization of its roots, is
\begin{equation}\label{eq:3-242}
BC - D = \frac{1+p}{p}\bigg[(2+p)\bigg[(1+x)^2 + T_1\,\frac{p^2}{(1+p)^2}\bigg] - R_1\,x\bigg].
\end{equation}

We may first observe that $D = 0$ and $BC - D = 0$ define the solutions which we have denoted by $R_1^{(c)}$ and $R_1^{(o)}$, respectively (see equations~\eqref{eq:3-227} and~\eqref{eq:3-228}). It will be recalled that these solutions give the Rayleigh numbers for the onset of stationary convection and overstability. This is consistent with equation~\eqref{eq:3-241}: for $\sigma_1 = 0$ is clearly a root of the equation if $D = 0$; while, if $\sigma_1$ should be purely imaginary, then on equating the real and imaginary parts of equation~\eqref{eq:3-241} separately, we should obtain $|\sigma_1|^2 = C = D/B$.

Now when $R_1 = 0$, $C$, $D$, and $BC - D$, are all positive; and $B$ is, of course, always positive. When $D > 0$, equation~\eqref{eq:3-241} clearly allows at least one real root $< 0$; denote this root by $-d$. When equation~\eqref{eq:3-241} allows a real negative root (such as $-d$), we can factorize it in the form
\begin{equation}\label{eq:3-243}
(\sigma_1 + d)(\sigma_1^2 + 2b\sigma_1 + c)(\sigma_1 + d) = 0.
\end{equation}

By comparison with equation~\eqref{eq:3-241}, we obtain the relations
\begin{equation}\label{eq:3-244}
B = 2b + d, \quad C = 2bd + c, \quad \text{and} \quad D = cd.
\end{equation}

From these relations it follows that
\begin{equation}\label{eq:3-245}
BC - D = 2b(2bd + c + d^2) = 2b(C + d^2);
\end{equation}
therefore,
\begin{equation}\label{eq:3-246}
b = \frac{BC - D}{2(C + d^2)} \quad \text{and} \quad c = \frac{D}{d}.
\end{equation}

As we have noted, when $R_1 = 0$, $BC - D > 0$ and $D > 0$; therefore
\begin{equation}\label{eq:3-247}
b > 0 \quad \text{and} \quad c > 0 \quad (\text{when } R_1 \to 0).
\end{equation}

The roots of equation~\eqref{eq:3-243}, besides $-d$, are
\begin{equation}\label{eq:3-248a}
\sigma_1 = -b \pm \sqrt{b^2 - c};
\end{equation}
by~\eqref{eq:3-247}, these two roots have also negative real parts. The state $R_1 = 0$ is, therefore, \emph{absolutely} stable since $\mathrm{re}(\sigma_1) < 0$ for all three roots.

Consider what happens when $R_1$ increases: $D$ and $BC - D$ decrease and one of them may become zero. It is evident that if a real root of~\eqref{eq:3-243} becomes zero, $D \to 0$, since either $c \to 0$ or $d \to 0$; while if the real part of a pair of complex roots becomes zero $b \to 0$, i.e.\ $BC - D \to 0$ (since $C > 0$, $C + d^2$ is necessarily positive). Therefore, as $R_1$ increases, \emph{either $D \to 0$ (in which case a real root tends to zero) or $BC - D \to 0$ (in which case the real part of a pair of complex roots tends to zero)}; and depending on whichever happens first, we shall have the onset of instability as stationary convection or as overstable oscillations. But this is exactly the premise on which our earlier discussion was based; it is now justified.

%% --- Section 30 ---
\section{On a method of discriminating the character of the marginal state. A general variational principle}\label{sec:3-30}

The discussion in the preceding section was restricted to the case of two free boundaries for which it was possible to obtain the required characteristic equation in an explicit form. We now return to the general case when this is not possible.

Replacing $\sigma$ by $i\sigma$ in equations~\eqref{eq:3-93}--\eqref{eq:3-95} and letting
\begin{equation}\label{eq:3-248}
F = (D^2 - a^2)(D^2 - a^2 - i\sigma)W - \bigg(\frac{2\Omega}{\nu}\,d^2\bigg) DZ,
\end{equation}
we have (cf.\ equations~\eqref{eq:3-105} and~\eqref{eq:3-106})
\begin{equation}\label{eq:3-249}
(D^2 - a^2 - i\sigma)Z = -\bigg(\frac{2\Omega}{\nu}\,d\bigg) DW
\end{equation}
and
\begin{equation}\label{eq:3-250}
(D^2 - a^2 - ip\sigma)F = -Ra^2 W.
\end{equation}

Equations~\eqref{eq:3-248} and~\eqref{eq:3-249} can be combined to give
\begin{equation}\label{eq:3-251}
[(D^2 - a^2 - i\sigma)^2(D^2 - a^2) + \mathscr{T} D^2] W = (D^2 - a^2 - i\sigma)F.
\end{equation}

Solutions of equations~\eqref{eq:3-248}--\eqref{eq:3-250} must be sought which satisfy the boundary conditions
\begin{equation}\label{eq:3-252}
\begin{aligned}
&W = F = 0 \quad \text{for} \quad z = 0 \;\text{and}\; 1, \\
\text{and}\quad &\textit{either}\; DW = 0 \;\text{and}\; Z = 0 \quad \text{(on a rigid surface),} \\
&\textit{or}\; D^2 W = 0 \;\text{and}\; DZ = 0 \quad \text{(on a free surface).}
\end{aligned}
\end{equation}

There are eight boundary conditions and the requirement that a solution of equations~\eqref{eq:3-248}--\eqref{eq:3-250} satisfying these conditions will determine for a given $a^2$ and $i\sigma$ (which can be real or complex) a sequence of possible values for $R$. These characteristic values of $R$ will in general be complex; and the requirement that $R$ be real implies a relation between the real and imaginary parts of $\sigma$. We are not, however, interested in the relationship which may exist between the real and imaginary parts of $\sigma$ and the other parameters of the problem. We are interested only in the marginal state and its characterization. From the discussion of the case of two free boundaries in \S~29, it is clear that it will suffice for our present purposes to determine the critical Rayleigh number for the onset of instability as stationary convection, and as overstability, for a given Taylor number; that which gives the lower Rayleigh number is the one that will be manifested first as the Rayleigh number is gradually increased, a disturbance in the horizontal plane characterized by the wave number $a$ (in the unit $1/d$) first becomes unstable by overstability when it reaches the value $R_0(a^2)$; and $\sigma_0(a^2)$ is the frequency (in the unit $\nu/d^2$) of the oscillations which are set up in the marginal state. To determine the critical Rayleigh number for the onset of overstability, we must determine the minimum of the function $R_0(a^2)$. We should then compare the minimum so obtained with the critical Rayleigh number for the onset of stationary convection; and that which is lower is the one that will prevail and be manifested at the onset of instability.

It should be apparent from the foregoing algorithm that an exact solution of the double characteristic value problem will be difficult if not impracticable. However, a variational method can be devised which makes the problem feasible.

\subsection*{(a) The variational principle}

Equations~\eqref{eq:3-248}--\eqref{eq:3-250} can be treated in the same manner as equations~\eqref{eq:3-104}--\eqref{eq:3-106} were treated in \S~26\,$(a)$. Thus, by multiplying equation~\eqref{eq:3-250}

\noindent by $F$ and integrating over the range of $z$, we obtain after an integration by parts the equation
\begin{equation}\label{eq:3-253}
\int_0^1 [(DF)^2 + (a^2 + ip\sigma)F^2]\,dz = Ra^2 \int_0^1 WF\,dz.
\end{equation}

The right-hand side of this equation requires us to consider
\begin{equation}\label{eq:3-254}
\int_0^1 WF\,dz = \int_0^1 W\bigg[(D^2 - a^2)^2 W - i\sigma(D^2 - a^2)W - \frac{2\Omega}{\nu}\,d^2\,DZ\bigg]\,dz.
\end{equation}

After several integrations by parts we find that
\begin{equation}\label{eq:3-255}
\int_0^1 WF\,dz = \int_0^1 \{[(D^2 - a^2)W]^2 + i\sigma[(DW)^2 + a^2 W^2]\}\,dz + \bigg(\frac{2\Omega}{\nu}\,d^2\bigg)\int_0^1 Z\,DW\,dz,
\end{equation}
the integrated parts always vanishing on account of the boundary conditions. Now making use of equation~\eqref{eq:3-249}, we find after further integrations by parts that
\begin{equation}\label{eq:3-256}
\int_0^1 WF\,dz = \int_0^1 \{[(D^2 - a^2)W]^2 + d^2 [(DZ)^2 + a^2 Z^2]\}\,dz + i\sigma \int_0^1 \{(DW)^2 + a^2 W^2 + d^2 Z^2\}\,dz.
\end{equation}

Thus, the result of multiplying equation~\eqref{eq:3-250} by $F$ and integrating over the range of $z$ is
\begin{equation}\label{eq:3-257}
\begin{split}
R &= \frac{\displaystyle\int_0^1 [(DF)^2 + a^2 F^2 + ip\sigma F^2]\,dz}{a^2\;\displaystyle\bigg\{\int_0^1 \{[(D^2 - a^2)W]^2 + d^2 [(DZ)^2 + a^2 Z^2]\}\,dz} \\
&\qquad\qquad + i\sigma \displaystyle\int_0^1 \{(DW)^2 + a^2 W^2 + d^2 Z^2\}\,dz\bigg\}} \\[6pt]
&= -\frac{I_1}{a^2 I_2} \;\text{(say)}.
\end{split}
\end{equation}

Now it can be shown (exactly as in \S~26\,$(a)$) that the variation $\delta R$ in $R$ given by equation~\eqref{eq:3-257} due to variations $\delta W$ and $\delta Z$ in $W$ and $Z$ compatible only with the boundary conditions on $W$, $Z$, and $F$ is given by
\begin{equation}\label{eq:3-258}
\delta R = -\frac{2}{a^2 I_2}\int_0^1 \delta F\{(D^2 - a^2 - ip\sigma)F + Ra^2 W\}\,dz.
\end{equation}

Consequently, \emph{if $\delta R = 0$ for all small arbitrary variations $\delta F$, then}
\begin{equation}\label{eq:3-259}
(D^2 - a^2 - ip\sigma)F = -Ra^2 W
\end{equation}
\emph{and conversely}. In other words, the characteristic values of $R$, for assigned $\sigma$, $a^2$, and $T$, have an extremal character; but unlike the case considered in \S~26\,$(a)$, they do not have a minimal character: indeed, they could not since they can be complex.

%% --- Section 31 ---
\section{The onset of convection as overstability: the solution for other boundary conditions}\label{sec:3-31}

The solution of the double characteristic value problem formulated in \S~30 can be obtained by the variational method by a procedure exactly analogous to that used in \S\S~27\,$(b)$ and~$(c)$. Thus, for obtaining the solution for the case of two rigid boundaries, we expand $F$ in a cosine series, as\footnote{The origin of $z$ has now been displaced so that its limits are $\pm\frac{1}{2}$.}
\begin{equation}\label{eq:3-260}
F = \sum_m A_m \cos\{(2m+1)\pi z\},
\end{equation}
and express $W$ and $Z$ as sums in the manner,
\begin{equation}\label{eq:3-261}
W = \sum A_m W_m \quad \text{and} \quad Z = \sum A_m Z_m.
\end{equation}

In view of equation~\eqref{eq:3-251}, $W_m$ is now a solution of the equation
\begin{equation}\label{eq:3-262}
[(D^2 - a^2 - i\sigma)^2(D^2 - a^2) + \mathscr{T} D^2] W_m = -c_{2m+1} \cos\{(2m+1)\pi z\},
\end{equation}
where
\begin{equation}\label{eq:3-263}
c_{2m+1} = (2m+1)^2\pi^2 + a^2 + i\sigma.
\end{equation}

The solution of equation~\eqref{eq:3-262} appropriate to the problem on hand is
\begin{equation}\label{eq:3-264}
W_m = c_{2m+1}\,\gamma_{2m+1}\cos\{(2m+1)\pi z\} + \sum_{j=1}^{3} B_j^{(m)} \cosh q_j\,z,
\end{equation}
where
\begin{equation}\label{eq:3-265}
\frac{1}{\gamma_{2m+1}} = c_{2m+1}^2[(2m+1)^2\pi^2 + a^2] + (2m+1)^2\pi^2\mathscr{T},
\end{equation}
the $B_j^{(m)}$ ($j = 1, 2, 3$) are constants of integration; and $q_j^2$ ($j = 1, 2, 3$) are the roots of the cubic equation,
\begin{equation}\label{eq:3-266}
(q^2 - a^2 - i\sigma)^2(q^2 - a^2) + \mathscr{T} q^2 = 0.
\end{equation}

The corresponding solution for $Z_m$ is
\begin{equation}\label{eq:3-267}
Z_m = -\frac{1}{\sqrt{\mathscr{T}}}\bigg\{(2m+1)\pi\gamma_{2m+1}\sin[(2m+1)\pi z] + \sum_{j=1}^{3}B_j^{(m)}\frac{q_j}{\chi_j}\sinh q_j\,z\bigg\},
\end{equation}
where
\begin{equation}\label{eq:3-268}
\chi_j = q_j^2 - a^2 - i\sigma \quad (j = 1, 2, 3).
\end{equation}

Comparing the solutions~\eqref{eq:3-264} and~\eqref{eq:3-267} with the solutions~\eqref{eq:3-146} and~\eqref{eq:3-161} obtained in \S~27\,$(a)$, we observe that they are identical except for

\noindent the different definitions of $q_j$, $\chi_j$, $c_{2m+1}$, and $\gamma_{2m+1}$. With these same redefinitions of the constants, the solutions for $B_j^{(m)}$ given in equations~\eqref{eq:3-154} and~\eqref{eq:3-155} apply equally well to the present case.

Substituting the solution for $W$ we have obtained in equation~\eqref{eq:3-250}, we now have
\begin{equation}\label{eq:3-269}
\begin{split}
&\sum A_m \bigg\{c_{2m+1}\gamma_{2m+1}\cos[(2m+1)\pi z] + \sum_{j=1}^3 B_j^{(m)}\cosh q_j\,z\bigg\} \\
&\qquad = Ra^2\sum A_m \bigg\{c_{2m+1}\gamma_{2m+1}\cos[(2m+1)\pi z] + \sum_{j=1}^3 B_j^{(m)}\cosh q_j\,z\bigg\}.
\end{split}
\end{equation}

Equation~\eqref{eq:3-269} is identical in form with equation~\eqref{eq:3-156}. The analysis following equation~\eqref{eq:3-156} up to, and inclusive of, equation~\eqref{eq:3-161} now applies.\footnote{Equations following equation~\eqref{eq:3-161} do not apply since these depend on the particular definitions of the roots $q_j$, etc.} In particular, we have the secular determinant:
\begin{equation}\label{eq:3-270}
\bigg\|\frac{1}{Ra^2}\,\frac{(2n+1)^2\pi^2 + a^2 + ip\sigma}{c_{2m+1}\gamma_{2m+1}}\,\delta_{mn} - \langle m|n\rangle\bigg\| = 0,
\end{equation}
where the matrix element $\langle m|n\rangle$ is given by equation~\eqref{eq:3-161}. The various quantities (such as $q_j$, $\chi_j$, $c_{2n+1}$, and $\gamma_{2m+1}$) appearing in equation~\eqref{eq:3-161} have the meanings now assigned to them.

The matrix $\langle m|n\rangle$ is symmetrical in $n$ and $m$. But the elements are complex, the symmetry does not ensure the reality of its characteristic roots; indeed they will in general be complex.

In the first approximation in which we retain only the $(0, 0)$-element of the secular matrix, the characteristic equation for $R$ becomes
\begin{equation}\label{eq:3-271}
R = \frac{1}{a^2}\,\frac{\pi^2 + a^2 + ip\sigma}{c_1\,\gamma_1\cdot\langle 0|0\rangle},
\end{equation}
where (cf.\ equation~\eqref{eq:3-161})
\begin{equation}\label{eq:3-272}
\begin{split}
\langle 0|0\rangle &= 2a^2_4 \Delta^{-1} \coth\tfrac{1}{2}q_1\bigg\{\tfrac{1}{2}(x_2 - x_1)(c_1 + x_3)^2 q_2\tanh\tfrac{1}{2}q_1 + {} \\
&\qquad + \tfrac{1}{2}(x_3 - x_2)(c_1 + x_2)^2 q_3\tanh\tfrac{1}{2}q_3 + \tfrac{1}{2}(x_3 - x_1)(c_1 + x_2)^2 q_2\tanh\tfrac{1}{2}q_2\bigg\}.
\end{split}
\end{equation}

It may be recalled that $\Delta$ in~\eqref{eq:3-272} is defined as in equation~\eqref{eq:3-155} but with the present meanings for the various quantities.

From our previous experience with variational methods, like the present, we may be confident that already the first approximation should give the required characteristic values to an accuracy of 1 or 2 per cent.

We shall now present the results of certain calculations based on equations~\eqref{eq:3-271} and~\eqref{eq:3-272}. The value
\begin{equation}\label{eq:3-273}
p = 0{\cdot}025
\end{equation}
was used in these calculations; this is approximately the value of $p$ for mercury at ordinary room temperatures. Some details regarding the method by which the $(R_c, T)$-relation was derived may be given.

For a chosen value of $a$, $R$ was evaluated in accordance with equations~\eqref{eq:3-271} and~\eqref{eq:3-272} for various assigned values of $\sigma$; and the value of $\sigma$ for which $R$ is real was deduced by interpolation. Thus for $T = 10^{10}$ and $a = 19{\cdot}8$ it was found that:
\begin{equation}\label{eq:3-274}
\begin{aligned}
R &= 1{\cdot}933 \times 10^6 - 5{\cdot}882 \times 10^4 i \quad \text{for } \sigma = 1{\cdot}420 \times 10^4; \\
R &= 1{\cdot}322 \times 10^6 + 9{\cdot}510 \times 10^4 i \quad \text{for } \sigma = 1{\cdot}500 \times 10^5; \\
R &= 1{\cdot}414 \times 10^6 + 0{\cdot}329 \times 10^4 i \quad \text{for } \sigma = 1{\cdot}489 \times 10^4.
\end{aligned}
\end{equation}

\begin{table}[ht]
\centering
\caption{Critical Rayleigh numbers and related constants for the onset of overstability in case both bounding surfaces are rigid and $p = 0{\cdot}025$}
\label{tab:XII}
\small
\begin{tabular}{ccccc}
\hline
$T$ & $a_c$ & $\sigma$ & $R_c$ & $p|\sigma|$ \\
\hline
$10^4$ & $3{\cdot}08$ & $4{\cdot}45 \times 10^1$ & $4{\cdot}39 \times 10^3$ & $0{\cdot}8692$ \\
$10^5$ & $4{\cdot}09$ & $5{\cdot}82 \times 10^2$ & $9{\cdot}31 \times 10^3$ & $1{\cdot}046$ \\
$5 \times 10^5$ & $8{\cdot}20$ & $2{\cdot}43 \times 10^3$ & $6{\cdot}39 \times 10^4$ & $0{\cdot}8862$ \\
$2 \times 10^6$ & $10{\cdot}28$ & $3{\cdot}90 \times 10^3$ & $1{\cdot}38 \times 10^5$ & $0{\cdot}5541$ \\
$10^7$ & $13{\cdot}46$ & $6{\cdot}81 \times 10^3$ & $3{\cdot}54 \times 10^5$ & $0{\cdot}4310$ \\
$10^{10}$ & $19{\cdot}7$ & $1{\cdot}29 \times 10^4$ & $1{\cdot}41 \times 10^6$ & $0{\cdot}2992$ \\
$10^{11}$ & $28{\cdot}75$ & $3{\cdot}47 \times 10^4$ & $5{\cdot}83 \times 10^6$ & $0{\cdot}2086$ \\
$10^{12}$ & $41{\cdot}7$ & $7{\cdot}18 \times 10^4$ & $2{\cdot}44 \times 10^7$ & $0{\cdot}1435$ \\
\hline
\end{tabular}
\end{table}

From these values, it can be estimated that
\begin{equation}\label{eq:3-275}
R = 1{\cdot}418 \times 10^6 \quad \text{for} \quad \sigma = 1{\cdot}4886 \times 10^4.
\end{equation}

By repeating such calculations for other $a$'s, the minimum of $R$ as a function of $a$ may be determined. Thus, in the example considered, it was found that:
\begin{equation}\label{eq:3-276}
\begin{aligned}
&\text{for } a = 19{\cdot}6, \quad \sigma = 1{\cdot}503 \times 10^4, \quad R = 1{\cdot}41765 \times 10^6; \\
&\text{for } a = 19{\cdot}7, \quad \sigma = 1{\cdot}496 \times 10^4, \quad R = 1{\cdot}4175 \times 10^6; \\
&\text{for } a = 19{\cdot}8, \quad \sigma = 1{\cdot}489 \times 10^4, \quad R = 1{\cdot}4177 \times 10^6.
\end{aligned}
\end{equation}

Consequently, it may be concluded that for $T = 10^{10}$, the critical Rayleigh number for the onset of overstability is $1{\cdot}4175 \times 10^6$ when $a = 19{\cdot}7$ and $\sigma = 1{\cdot}496 \times 10^4$.

In Table~XII the results of such calculations are summarized. Similar, but less extensive, calculations were undertaken for the case when one of the bounding surfaces is free and the other is rigid. The results of these calculations are given in Table~XIII.

\begin{table}[ht]
\centering
\caption{Critical Rayleigh numbers and related constants for the onset of overstability in case one bounding surface is free and the other is rigid and $p = 0{\cdot}025$}
\label{tab:XIII}
\small
\begin{tabular}{ccccc}
\hline
$T$ & $a_c$ & $\sigma$ & $R_c$ & $p|\sigma|$ \\
\hline
$10^5$ & $5{\cdot}85$ & $1{\cdot}44 \times 10^3$ & $1{\cdot}71 \times 10^4$ & $0{\cdot}9125$ \\
$3 \times 10^5$ & $15{\cdot}38$ & $1{\cdot}64 \times 10^4$ & $4{\cdot}81 \times 10^4$ & $0{\cdot}3601$ \\
$10^6$ & $0{\cdot}5$ & $7{\cdot}46 \times 10^4$ & $4{\cdot}87 \times 10^5$ & $0{\cdot}1492$ \\
\hline
\end{tabular}
\end{table}

The results on the critical Rayleigh number, for the onset of instability as stationary convection and as overstable oscillations for all three sets of boundary conditions, are all included in Figs.~21 and~22 (p.~96).

%% --- Section 32 ---
\section{The case $\mathfrak{p} = 0$}\label{sec:3-32}

The case when the Prandtl number $p$ tends to zero is singular with respect to the onset of overstable oscillations: thus, the asymptotic relations~\eqref{eq:3-234}, \eqref{eq:3-235}, and~\eqref{eq:3-237} are valid for $p^2 T \to \infty$, and they cannot be used when $p = 0$. The case $p = 0$ should be treated separately.

Returning then to equation~\eqref{eq:3-222} and setting $p = 0$, we have\footnote{We are restricting our present considerations to the case of two free boundaries.}
\begin{equation}\label{eq:3-277}
R_c^{(o)} = 2\pi^4\,\frac{(1+x)^3}{x};
\end{equation}
apart from the factor 2, this is the same formula which obtains in the absence of rotation (cf.\ equation~II\;\eqref{eq:3-193}). The minimum Rayleigh number occurs for $x = \frac{1}{2}$ when
\begin{equation}\label{eq:3-278}
R_c^{(o)} = 13{\cdot}5\pi^4 = 1315.
\end{equation}

The frequency of the overstable oscillations, when instability occurs, can be obtained from equation~\eqref{eq:3-221} by putting $p = 0$ and $x = \frac{1}{2}$; we get
\begin{equation}\label{eq:3-279}
\sigma_1^2 = \tfrac{4}{9}T_1 - 2{\cdot}25.
\end{equation}

Ignoring $2{\cdot}25$ in this equation and making use of the relation~\eqref{eq:3-236}, we have
\begin{equation}\label{eq:3-280}
p = (\tfrac{2}{3})^{\frac{1}{2}} 2\Omega,
\end{equation}
where $p$ is now the (circular) frequency in seconds.

The frequency given by equation~\eqref{eq:3-280} is different from the natural frequency of the rotating fluid only by the factor $(2/3)^{\frac{1}{2}}$. It would, therefore, be physically correct to say that \emph{in the limit $p = 0$, instability sets in by exciting the natural modes of oscillation of the rotating fluid}.

%% --- Section 33 ---
\section{Thermodynamic significance of the variational principles}\label{sec:3-33}

We shall now show that as in the case of the simple B\'enard problem (\S~14), the variational principles established in \S\S~26\,$(a)$ and~30 have similar thermodynamic meanings.

Consider the average rate of viscous dissipation of energy in a unit column of the fluid and the average rate at which the buoyancy force, $g\alpha\Theta$ ($= g\alpha\beta$) releases energy in the same unit column. These are given by (equations~II~\eqref{eq:3-175} and~\eqref{eq:3-178})
\begin{equation}\label{eq:3-281}
\epsilon_v = -\frac{p\nu}{d^2}\int_0^1\big\{\langle u(D^2 - a^2)u\rangle + \langle u(D^2 - a^2)u\rangle + \langle v(D^2 - a^2)v\rangle\big\}\,dz
\end{equation}
and
\begin{equation}\label{eq:3-282}
\epsilon_g = \rho g\alpha\int_0^1 \langle\delta w\rangle\,dz,
\end{equation}
where angular brackets signify that the quantity enclosed is averaged over the horizontal plane.

Without loss of generality for our present purposes, we may suppose that the expressions for $u$, $v$, and $w$ are those given in equations~\eqref{eq:3-186} and~\eqref{eq:3-187}. According to these equations, we have, for example,
\begin{equation}\label{eq:3-283}
\int_0^1 \langle u(D^2 - a^2)u\rangle\,dz = -\frac{1}{4a^4}\int_0^1 \{a_x^2\,DW(D^2 - a^2)\,DW + a_y^2 d^2 Z(D^2 - a^2)Z\}\,dz.
\end{equation}

It should be noted that the cross terms involving the products
\[
DW(D^2-a^2)Z \quad \text{and} \quad Z(D^2-a^2)W
\]
do not make any contributions to the foregoing expression; this is due to the quite general circumstance that the waves in the horizontal plane associated with $DW$ and $Z$ in the solution for $\mathbf{u}_\perp$ are exactly out of phase.

We have an expression similar to~\eqref{eq:3-283} for the term in $v$ in equation~\eqref{eq:3-281}; also,
\begin{equation}\label{eq:3-284}
\int_0^1 \langle w(D^2 - a^2)w\rangle\,dz = \pm\tfrac{1}{4}\int_0^1 W(D^2 - a^2)W\,dz.
\end{equation}

Combining these results, we have (cf.\ equation~II~\eqref{eq:3-178})
\begin{equation}\label{eq:3-285}
\epsilon_v = -\frac{p\nu}{4d^2 a^2}\int_0^1 \{a^2 W(D^2 - a^2)W + DW(D^2 - a^2)DW + d^2 Z(D^2 - a^2)Z\}\,dz.
\end{equation}

After integrating by parts the second and the third terms on the right-hand side, we find
\begin{equation}\label{eq:3-286}
\epsilon_v = \frac{p\nu}{4d^2 a^2}\int_0^1 \{[(D^2 - a^2)W]^2 + d^2 [(DZ)^2 + a^2 Z^2]\}\,dz.
\end{equation}

In the further reduction of the expression for $\epsilon_g$, we must distinguish the stationary and the oscillatory cases.

\subsection*{(a) The case when the marginal state is stationary}

In this case, the equation relating $\Theta$ and $w$ is (cf.\ equation~\eqref{eq:3-86})
\begin{equation}\label{eq:3-287}
w = -\frac{\kappa}{\beta d^2}(D^2 - a^2)\Theta.
\end{equation}

Accordingly,
\begin{equation}\label{eq:3-288}
\epsilon_g = -\frac{\rho g\alpha\kappa}{4\beta d^2}\int_0^1 \langle \Theta(D^2 - a^2)\Theta\rangle\,dz = \frac{\rho g\alpha\kappa}{4\beta d^2}\int_0^1 \Theta(D^2 - a^2)\Theta\,dz.
\end{equation}

After an integration by parts, we have
\begin{equation}\label{eq:3-289}
\epsilon_g = \frac{\rho g\alpha\kappa}{4\beta d^2}\int_0^1 [(D\Theta)^2 + a^2\Theta^2]\,dz.
\end{equation}

On the other hand, according to equations~\eqref{eq:3-98} and~\eqref{eq:3-104},
\begin{equation}\label{eq:3-290}
\Theta = -\frac{\nu}{g\alpha d^2 a^2}\,F.
\end{equation}

In terms of $F$ the expression for $\epsilon_g$ becomes
\begin{equation}\label{eq:3-291}
\epsilon_g = \frac{p\kappa^2}{4g\alpha\beta d^4 a^2}\int_0^1 [(DF)^2 + a^2 F^2]\,dz.
\end{equation}

In a steady state, the kinetic energy dissipated by viscosity must be balanced by the internal energy released by the buoyancy force. We must, therefore, require
\begin{equation}\label{eq:3-292}
\epsilon_v = \epsilon_g.
\end{equation}

Equations~\eqref{eq:3-286} and~\eqref{eq:3-291} now give
\begin{equation}\label{eq:3-293}
R = \frac{g\alpha\beta d^4}{\kappa\nu} = \frac{\displaystyle\int_0^1 [(DF)^2 + a^2 F^2]\,dz}{a^2\;\displaystyle\int_0^1 \{[(D^2 - a^2)W]^2 + d^2 [(DZ)^2 + a^2 Z^2]\}\,dz},
\end{equation}

\noindent But this is exactly the same expression for $R$ which, according to the variational principle of \S~26\,$(a)$, attains its absolute minimum for the marginal state.

\subsection*{(b) The case when the marginal state is oscillatory}

The question now arises: how should one modify under circumstances of overstability, the principle which equates the rate of irreversible dissipation of energy with the rate of liberation of the thermodynamically available energy by the buoyancy force, as a criterion for marginal stability? It would appear natural that we generalize the principle along the following lines.

If the motion is periodic in time, then the kinetic energy of the fluid, as well as the release of the potential energy by the buoyancy force, will be subject to similar variations. Suppose that all quantities which describe the perturbation vary with a circular frequency $p$ so that all the amplitudes have a time-dependent factor $e^{ipt}$. Then
\begin{equation}\label{eq:3-294}
\mathbf{u}_i\frac{\partial \mathbf{u}_i}{\partial t} = ip\mathbf{u}_i^2.
\end{equation}

We must allow for this change in the kinetic energy (per unit mass) in writing an equation of energy balance. Thus, it would appear that the quantity we must equate to $\epsilon_v$ is
\begin{equation}\label{eq:3-295}
\epsilon_v + ipp\int_0^1 \langle\mathbf{u}_i^2\rangle\,dz.
\end{equation}

Reverting to $d$ as the unit of distance and defining $\sigma$ ($= pd^2/\nu$) as in equation~\eqref{eq:3-92}, we have to consider
\begin{equation}\label{eq:3-296}
\epsilon_v + i\sigma\frac{p\nu}{d^2}\int_0^1 \langle\mathbf{u}_i^2\rangle\,dz
\end{equation}
in place of $\epsilon_v$ for the stationary case. The criterion for the occurrence of an oscillatory marginal state should then be
\begin{equation}\label{eq:3-297}
\epsilon_v + i\sigma\frac{p\nu}{d^2}\int_0^1 \langle\mathbf{u}_i^2\rangle\,dz = \epsilon_g = g\alpha\rho\int_0^1 \langle\delta w\rangle\,dz,
\end{equation}
where in evaluating $\epsilon_g$ we must allow for the fact that $\Theta$ and $w$ are no longer stationary.

At first sight, the appearance of the imaginary $i$, and complex numbers generally, in the equation for the energy balance, may strike one as very odd. Its origin must be traced to the fact that the oscillations in the velocity and in the acceleration are out of phase. Consequently, the

\noindent excess (or defect) of energy dissipated during one phase of the cycle must be exactly compensated by a similar excess (or defect) of energy liberated in a synchronous manner. It is this need for synchronism which determines the period of oscillation, as well as the Rayleigh number, as characteristics of the marginal state. The mathematical counterpart of this is that the determination of $R$ and $\sigma$ is through the solution of a double characteristic value problem.

We shall now verify that the thermodynamic principle, as reformulated above, is in agreement with the variational principle of \S~30. The expression~\eqref{eq:3-286} for $\epsilon_v$ continues to be valid. And we must now evaluate the additional term in $\langle\mathbf{u}_i^2\rangle$ in equation~\eqref{eq:3-297}. Considering the contribution to the integral by the $x$-component of the velocity, we have (cf.\ equation~\eqref{eq:3-187})
\begin{equation}\label{eq:3-298}
\int_0^1 \langle u^2\rangle\,dz = \frac{1}{4a^4}\int_0^1 [a_x^2(DW)^2 + a_y^2 d^2 Z^2]\,dz.
\end{equation}

We have a similar contribution from the $y$-component of the velocity; and altogether we have
\begin{equation}\label{eq:3-299}
i\sigma\frac{p\nu}{d^2}\int_0^1 \langle\mathbf{u}_i^2\rangle\,dz = i\sigma\,\frac{p\nu}{4d^2 a^2}\int_0^1 [a^2 W^2 + a^2(DW)^2 + a^2 d^2 Z^2]\,dz.
\end{equation}

This must be added to the expression~\eqref{eq:3-286} giving $\epsilon_v$.

Turning to the evaluation of $\epsilon_g$, we must now use the relation (see equation~\eqref{eq:3-93})
\begin{equation}\label{eq:3-300}
(D^2 - a^2 - ip\sigma)\Theta = -\bigg(\frac{\beta\,d^2}{\kappa}\bigg)w
\end{equation}
for eliminating $w$. Accordingly, we now have
\begin{equation}\label{eq:3-301}
\epsilon_g = -\frac{g\alpha\rho\kappa}{4\beta d^2}\int_0^1 \Theta(D^2 - a^2 - ip\sigma)\Theta^*\,dz = \frac{g\alpha\rho\kappa}{4\beta d^2}\int_0^1 [(D\Theta)^2 + (a^2 + ip\sigma)\Theta^2]\,dz.
\end{equation}

The relation between $\Theta$ and $F$ is unaltered and is given, as before, by~\eqref{eq:3-290}. Thus,
\begin{equation}\label{eq:3-302}
\epsilon_g = \frac{p\kappa^2}{4g\alpha\beta d^4 a^2}\int_0^1 [(DF)^2 + (a^2 + ip\sigma)F^2]\,dz.
\end{equation}

Combining equations~\eqref{eq:3-286}, \eqref{eq:3-297}, \eqref{eq:3-299}, and~\eqref{eq:3-302}, we obtain
\begin{equation}\label{eq:3-303}
R = \frac{\displaystyle\int_0^1 [(DF)^2 + (a^2 + ip\sigma)F^2]\,dz}{a^2\;\displaystyle\bigg\{\int_0^1\{[(D^2 - a^2)W]^2 + d^2[(DZ)^2 + a^2 Z^2]\} + i\sigma[(DW)^2 + a^2 W^2 + d^2 Z^2]\}\,dz\bigg\}}
\end{equation}
and this is, indeed, the expression for $R$ which formed the basis of the variational principle of \S~30.

We are thus enabled to formulate the following general principle.

\emph{Thermal instability as stationary convection will set in at the minimum (adverse) temperature gradient which is necessary to maintain a balance between the rate of dissipation of energy by viscosity and the rate of liberation of the thermodynamically available energy by the buoyancy force acting on the fluid. Likewise, the onset of thermal instability will be as overstable oscillations if it is possible (at a lower adverse temperature gradient) to balance in a synchronous manner the periodically varying amounts of kinetic energy with similarly varying amounts of dissipation and liberation of energy.}

%% --- Section 34 ---
\section{The case when $\boldsymbol{\Omega}$ and $\mathbf{g}$ act in different directions}\label{sec:3-34}

We shall now briefly consider the case when $\boldsymbol{\Omega}$ and $\mathbf{g}$ act in different directions. For this purpose we must return to equations~\eqref{eq:3-84} and~\eqref{eq:3-85} in which the assumption that $\boldsymbol{\Omega}$ and $\mathbf{g}$ are parallel has not yet been made.

Let $\boldsymbol{\Omega}$ be inclined at an angle $\vartheta$ to the direction of the vertical. Also, let the direction of the $x$-axis be so chosen that $\boldsymbol{\Omega}$ lies in the $xz$-plane. Then
\begin{equation}\label{eq:3-304}
\boldsymbol{\lambda} = (0, 0, 1) \quad \text{and} \quad \boldsymbol{\Omega} = \Omega(\sin\vartheta, 0, \cos\vartheta),
\end{equation}
and equations~\eqref{eq:3-84} and~\eqref{eq:3-85} become
\begin{equation}\label{eq:3-305}
\frac{\partial}{\partial t}\nabla^2 w = \nu\nabla^4 w + 2\Omega\bigg(\cos\vartheta\,\frac{\partial \zeta}{\partial z} + \sin\vartheta\,\frac{\partial \zeta}{\partial x}\bigg) + g\alpha\nabla_1^2\Theta
\end{equation}
and
\begin{equation}\label{eq:3-306}
\frac{\partial}{\partial t}\nabla^2 w = g\alpha\bigg(\frac{\partial^2}{\partial x^2} + \frac{\partial^2}{\partial y^2}\bigg)\Theta + \nu\nabla^4 w - 2\Omega\bigg(\cos\vartheta\,\frac{\partial^2}{\partial z^2} + \sin\vartheta\,\frac{\partial}{\partial x}\frac{\partial}{\partial z}\bigg)\zeta.
\end{equation}

The remaining equation~\eqref{eq:3-86} is unaffected by this generalization.

If we seek solutions of equations~\eqref{eq:3-86}, \eqref{eq:3-305}, and~\eqref{eq:3-306} which are independent of $x$ and are of the forms
\begin{equation}\label{eq:3-307}
w = W(z)\cos\frac{a_y}{d}\,y, \quad \zeta = Z(z)\cos\frac{a_y}{d}\,y, \quad \text{and} \quad \Theta = \Theta(z)\cos\frac{a_y}{d}\,y,
\end{equation}
then equations~\eqref{eq:3-305} and~\eqref{eq:3-306} reduce to equations~\eqref{eq:3-87} and~\eqref{eq:3-88} with the only difference that $\Omega$ is replaced by $\Omega\cos\vartheta$. Consequently, if we restrict ourselves to the onset of instability as rolls in the $x$-direction, then all of the discussion of the preceding sections, with the exception of those parts of \S~28 which specifically deal with general cell patterns, will apply if we interpret $\Omega$ everywhere to mean the component of $\boldsymbol{\Omega}$ in the direction of $\mathbf{g}$.

The question, whether lower Rayleigh numbers can be attained for the

\noindent marginal states by considering more general patterns of motion than rolls, can be answered only by solving the necessary characteristic value problems in which $a_x$ and $a_y$ are explicitly retained; and minimizing the lowest characteristic values, for given $a_x$ and $a_y$, as functions of these two parameters. Such an analysis has not been carried out for the present problem; but it has been carried out for the closely related problem of the inhibition of thermal convection by an impressed magnetic field (Chapter~IV, \S~47). If the conclusions reached in that context can be extended to the present problem, then it would appear that the minimum Rayleigh numbers do indeed occur for the onset of instability as rolls. This apparently paradoxical conclusion is clarified in Chapter~IV, \S~47.

%% --- Section 35 ---
\section{Experiments on the onset of thermal instability in rotating fluids}\label{sec:3-35}

% PLACEHOLDER: Section 35 content

%% --- Bibliographical Notes ---
\section*{BIBLIOGRAPHICAL NOTES}
\addcontentsline{toc}{section}{Bibliographical Notes}

% PLACEHOLDER: Bibliographical Notes content
